%%%%%%%%% JPS abstract %%%%%%%%%%%%%%%%%%%%%%%%%%%%%%%%%%%%%%%%%%
\documentclass[12pt,a4paper,upLaTeX]{jsarticle}

%%%%%%%%% packages %%%%%%%%%%%%%%%%%%%%%%%%%%%%%%%%%%%%%%%%%%%%%%
\usepackage[dvipdfmx]{graphicx} % Include figure files
\usepackage[%                           % 余白の設定
%mag=1400,%                              jarticle の場合（14ptに）
dvipdfm,truedimen,%
top=30truemm,bottom=50truemm,%
left=20truemm,right=20truemm]{geometry}
\usepackage{fancyhdr}
\usepackage{amsmath,amssymb}
\usepackage{bm}
\pagestyle{fancy}

%%%%%%%%% header %%%%%%%%%%%%%%%%%%%%%%%%%%%%%%%%%%%%%%%%%%%%%%%%
\begin{document}

\fancyhf{}
\setlength{\headheight}{60pt}   % ある程度大きめに
\fancyhead[R]{大阪大学大学院 理学研究科 物理学専攻 博士前期課程2年\\
% 学籍番号 24B25011\\
加島 駿一\\}

\setlength{\footskip}{20pt}
% \fancyhf{}
\fancyfoot[C]{JPS2026Autumn Summary}

\fancyfoot[R]{\thepage}

\vspace{50pt}
\begin{center}
{\gt \Large JPS2026Autumn Summary }\\[14pt]

\end{center}

\vspace{10pt}
%%%%%%%%%%%%%%%%%%%%%%%%%%%%%%%%%%%%%%%%%%%%%%%%%%

\subsubsection*{p.1}
\begin{itemize}
\item 研究の背景について説明します
\item チャームクォークを含むバリオンである、チャームバリオンではクォークの質量差によって、内部に存在する軽いダイクォーク相関が浮き彫りになります。つまりバリオン内部のクォークの運動を見ることができると言えます
\item そしてこれはチャームバリオンの励起準位構造や励起状態の生成率・崩壊分岐比に現れるので、我々はJ-PARCにてチャームバリオンの分光実験E50を計画しています
\item 運動量20 GeV/cのpionビームを液体水素標的に衝突させ、生成されるD*-崩壊由来のK+.pi-,pi-を前方のスペクトロメーターで捉えることで質量欠損法でチャームバリオンの質量スペクトルを測ります
\item この反応断面積は数ナノbと見積もられており、非常にレアな反応です。十分な統計量を得るために30 mega cpsの大強度のビームを用います
\item 我々はこの実験においてトリガーレスの連続読み出しでデータ収集を行い、この非常に高い反応レートで、ハードウェアトリガーでトリガーするのではなく全てソフトウェアでフィルタリングします
\end{itemize}

\subsubsection*{p.2}
\begin{itemize}
\item 我々のDAQシステムではハードウェアレベルでのトリガーは存在せず、全てコンピュータに取り込んだ上で必要なデータを必要なだけとるようなフィルタリングを行います
\item NestDAQという連続読み出し式DAQフレームワーク上で、MARQスペクトロメーターのためのオンラインフィルターを実装し、開発・性能評価が進んでいます
\item 右のイメージで示しているように、ビーム粒子の90%はターゲットと反応せずに通り抜けてしまう、つまり興味のないイベントが多いという背景があるので、オンラインで散乱粒子の有無を判別することでビームスルーイベントを削減することで、賢くデータを貯めることがオンラインフィルターの当面の目標です
\item そのためにドリフトチェンバーやシンチレーションファイバートラッカーを用いてトラッキングを行うようなオンラインフィルターの開発いました。本公演のテーマはその処理性能評価です
\end{itemize}

\subsubsection*{p.3}
\begin{itemize}
\item DAQフレームワークであるNestDAQについて説明します
\item この絵のピンクの枠で囲んだ部分がNestDAQなのですが、フロントエンドの読み出し基盤から送られてきたデータを受け取ってある時間幅でブロック化する部分や、各フロントエンドのデータをひとまとめに組み立てる部分、次に検出器のヒットのコインシデンスを探してイベントパケットを切り出す部分、次にオンラインフィルターがイベントを取捨選択し、ディスクに書き込むという構成要素で動きます
\item 赤枠で囲った部分は入力データレートに応じてスケールすることができます
\begin{itemize}
\item 例えば、オンラインフィルターの負荷が高いから並列で動作させる数をこんな感じで増やすとか
\item 必要な並列プロセス数は処理すべきデータレートと1プロセスあたりの処理できるデーターレートの比になります
\end{itemize}
\item システムにスケーラビリティがあって理論上無限のスループットがあるとはいえ、計算資源には限りがあり、どれくらいの計算資源でどんな処理をどんなスピードでできるのかを見積もることは非常に重要です
\item 具体的にはE50実験でビームスルーイベントを削減するにはどれくらいの計算資源が必要かを見積もるために、過去のDAQのテストベンチで取得したデータからドリフトチェンバーのトラッキングを行うフィルター1プロセスあたりの処理性能評価を行ったのでそれについて報告します
\end{itemize}

\subsubsection*{p.4}
\begin{itemize}
\item 性能評価の測定対象について説明します
\item 右の絵がオンラインフィルターの処理内容なのですが、赤矢印で示した、他のデバイスとのデータのやり取りなどを含めたオンラインフィルターの動作全体にかかる時間とトラッキングや条件判定を行う関数にかかる時間を測定します
\begin{itemize}
\item ちなみに、NestDAQユーザーがフィルタリングしようと思うと基本的に青色矢印の部分の処理を適宜記述することになります
\end{itemize}
\item 処理したデータ量を測定した時間で割ることで処理できるデータレートが計算できます
\item 本発表の測定では右の表の通りの計算機で測定を行いました
\end{itemize}

\subsubsection*{p.5}
\begin{itemize}
\item 測定に用いたデータについて説明します
\item 2024年4月にJ-PARC K1.8BRビームラインにて行われた連続読み出し式DAQのテストで取得したデータを用います
\item 写真の通り複数の検出器で構成されており、青色の文字で示したタイミングカウンタと赤色の文字で示したトラッカーの合計チャンネル数は4300chほどです
\item このテストベンチに1 GeV/c, 0.95 Mcpsのビームを照射し、フロントエンドからは260 MB/sほどのデータフローがありました
\item このようなデータに対して、右の図でいうところのビーム上流からUTOFというプラスチックシンチレータ、KLDCというドリフトチェンバー、LTOFというプラスチックシンチレータに着目し、
\item UTOFとLTOFのヒットのコインシデンスを基準に±1000 nsのウィンドウでイベントをスライスし、それをオンラインフィルターに解析させて取捨選択させます
\item このテストベンチの詳細に関してはこの論文をご参照ください
\end{itemize}

\subsubsection*{p.6}
\begin{itemize}
\item 今回使用したドリフトチェンバーは±45度傾いたU,U’,V,V'の面を持ち、それぞれが128本のワイヤを持ちます。セルの構造は六角形です。
\item また、例えばUとU'のレイヤーはセルの半分だけずらして配置してあります
\item このようなドリフトチェンバーのヒットに対して次のような処理を行います
\item まず最初に生データからドリフトチェンバーのヒット情報へと変換を行います
\begin{itemize}
\item その際に、自作の汎用チャンネルマップを用いて読み出し基盤のIPアドレスや読み出しチャンネルの情報からそれがどんな検出器のヒットなのかという情報や、キャリブレーションパラメータを拾います
\item このオンラインフィルターの性能評価はそのチャンネルマップの統合試験も兼ねて行っています
\item これらの情報からTDCをドリフト時間やドリフト距離へと変換し、
\end{itemize}
\item 次に各ペアになっている面についてヒットのクラスタリングを行い、
\item 最後にトラッキングを行います
\begin{itemize}
\item トラッキングには、直線トラックとドリフトチェンバーのヒット位置の残差について最小二乗法フィッティングを行い、直線の切片と傾きを求めます
\end{itemize}
\item こうして求めたトラックの情報を用いてビームと散乱粒子を判別できるようなフィルタープロセスになっています
\end{itemize}

\subsubsection*{p.7}
\begin{itemize}
\item 処理性能の測定結果について話します
\item これは前半のスライドにあった絵ですが、赤矢印で示したオンラインフィルター全体の処理にかかる時間から、オンラインフィルターの処理できるデータレートを測定しました
\item 横軸にスループットをとっており、平均して123.4 MB/sで処理できるとわかりました
\begin{itemize}
\item この数字がどうだってのは次のページで話します
\end{itemize}
\item また、この赤矢印と青矢印の時間は平均1:0.9でした
\begin{itemize}
\item つまりオンラインフィルターの処理の中でほとんどが本質的な条件判定にかかる時間であるとわかります
\end{itemize}
\item また、青矢印の処理のうち、生データからヒット情報への変換、クラスタリング、フィッティングそれぞれにかかる時間の割合についても測定を行い、下のグラフのようになりました。横軸は割合です
\item これを見ると大部分がフィッティングの計算にかかる時間でした
\end{itemize}

\subsubsection*{p.8}
\begin{itemize}
\item 前のページで測定した処理できるデータレートを用いて、処理に必要な計算資源を見積もります
\item 今回のデータは、2sのビーム取り出しと2.2秒のビームが来ない時間のサイクル構造を持ったビームで取ったデータなので、必要プロセス数の計算にはビームのduty比をかける必要があります
\item オンサイトでフィルタリングをしようと思うと、必要なプロセス数は処理すべきデータレートと処理できるデータレートがこうなので、2.02個、つまりCPUが3スレッドあれば余裕で処理できる性能であることがわかりました
\item E50本番の想定データレートで換算すると少なくとも59個のトラッキングフィルタープロセスが要求されます
\begin{itemize}
\item 実現可能性としては十分にあると考えています
\end{itemize}
\end{itemize}

\subsubsection*{p.9}
まとめ

\subsubsection*{質疑応答}
\begin{itemize}
\item フィッティングまでやってるけど、クラスタリングだけ or クラスターのジオメトリ or クラスター特徴量を抜き出す云々だけで1ケタの削減は可能なのでは?
\begin{itemize}
\item クラスタリングだけでは無理だけど、たしかにそうかも
\end{itemize}
\item フィッティングにかかる時間のもっと詳細は?
\begin{itemize}
\item 計測してないです
\end{itemize}
\item 計算機のスペックが想像がわかない
\begin{itemize}
\item コンシューマの最上位付近
\end{itemize}
\item スループットの幅広いよね、その幅の広さの要因は?また、スループットが低いものを想定して安全ファクターをかけなきゃいけないのじゃない?
\begin{itemize}
\item 幅の広さの要因はわからないけど、スループットが低いものは処理時間が長いものが多かった
\item 安全ファクターは不要だと思ってる。連続読み出しはソフトウェアのバッファさえ埋まらなければ大丈夫だから、処理すべきデータレートは長い時間で計算すればok。1ランあたりのデータ量/1ランあたりの処理時間で算出している。
\end{itemize}
\end{itemize}


\end{document}